\title{FlashAttention-3: Fast and Accurate Attention with Asynchrony and Low-precision}

\begin{abstract}
Attention serves as a foundational component in the widespread Transformer framework, yet it remains a principal computational bottleneck, particularly for large language models and scenarios involving extended contexts. The \textsc{FlashAttention} method proposed significant GPU acceleration of attention by reducing memory accesses. Nonetheless, this approach does not fully exploit the advanced functionalities of newer hardware, as \textsc{FlashAttention-2} manages to utilize just 35\% of the H100 GPU’s capabilities. In this work, we propose three novel strategies to enhance attention performance on Hopper GPUs: (1) leveraging the asynchronous operation of Tensor Cores and TMA to overlap computation with data transfer via warp-specialization, (2) introducing an interleaved execution of block-wise matrix multiplication and softmax routines, and (3) employing block quantization together with incoherent processing, capitalizing on native hardware support for low-precision FP8 calculations. Through these innovations, our approach, \textsc{FlashAttention-3}, delivers a 1.5-2.0$\times$ increase in speed on H100 GPUs—FP16 computations reach up to 740 TFLOPs/s, corresponding to 75\% hardware utilization, while FP8 calculations approach 1.2 PFLOPs/s. Furthermore, our evaluation shows that FP8 \textsc{FlashAttention-3} reduces numerical error by a factor of 2.6$\times$ compared to standard FP8 attention baselines.
\end{abstract}

\section{Introduction}
\label{sec:intro}

Within the Transformer architecture~\citep{vaswani2017attention}, the attention mechanism represents the central computational constraint, as the calculation of self-attention between queries and keys scales quadratically with respect to the sequence length. Expanding the capacity of attention mechanisms to manage extended contexts can facilitate novel capabilities, such as modeling and performing reasoning across multiple lengthy documents~\citep{guo2021longt5,shaham2022scrolls,peng2023yarn}, handling extensive code repositories~\citep{roziere2023code, li2023starcoder}, processing different modalities like high-resolution imagery~\citep{chen2022scaling}, audio~\citep{gulati2020conformer}, and video~\citep{ho2022video}, as well as supporting new practical applications, including user interfaces with extensive historical interactions~\citep{sun2019bert4rec} and agent workflows spanning prolonged timeframes~\citep{yao2022react}. Consequently, there has been substantial research attention directed toward accelerating attention mechanisms for longer input sequences, either through approximation methods~\citep{katharopoulos2020transformers,choromanski2020rethinking,
tay2020efficient}, advances in software implementation (\citep{rabe2021self,dao2022flashattention,kwon2023efficient}), or by developing alternative model architectures~\citep{peng2023rwkv,sun2023retentive,gu2023mamba}.

In this paper, we extend the foundation established by \citet{dao2022flashattention}, who designed exact-attention algorithms that explicitly leverage the GPU execution model and hardware attributes at a high level. Dao et al. proposed \textsc{FlashAttention} in \citep{dao2022flashattention}, introducing an innovative tiling approach that merges all attention operations into one GPU kernel, thereby avoiding the need for intermediate data transfers to slower global memory.

Subsequently, \citet{dao2023flashattention2} refined this approach with \textsc{FlashAttention-2}, restructuring the algorithm to distribute computations across the sequence length dimension and to process the key and value matrices in block-wise fashion during the forward pass's inner loop. This modification aimed to enhance GPU workload balancing and occupancy. Nevertheless, our analysis reveals that \textsc{FlashAttention-2} still demonstrates suboptimal hardware usage on the latest GPUs compared with highly optimized GEMM (matrix multiplication) kernels—for example, achieving only 35\% of peak utilization on the Hopper H100 GPU as opposed to the 80–90\% reached by these specialized kernels. Some of this efficiency gap stems from implementation-level limitations, such as reliance on Ampere-targeted instructions instead of those optimized for Hopper Tensor Cores. Recent studies, including ThunkerKitten~\citep{spector2024thunder} and the release of cuDNN 9~\citep{cudnn9}, indicate that incorporating Hopper-specific instructions and tile-centric abstractions can both accelerate attention computations and simplify the overall implementation.

At a foundational level, the algorithm underlying \textsc{FlashAttention-2} operates within a basic synchronous framework, deliberately omitting explicit mechanisms for exploiting hardware asynchrony or low-precision arithmetic. The emergence of asynchrony is driven by hardware optimizations tailored to expedite critical components of machine learning workloads: specialized hardware modules manage matrix multiplications (via Tensor Cores) and facilitate efficient memory transfer (through the Tensor Memory Accelerator, or TMA), while a separate set of CUDA cores is tasked with general computation—including logic, integer operations, and floating-point calculations. Progress in hardware-supported low-precision formats, as exemplified by the introduction of FP8 in Hopper and FP4 in Blackwell (preceded by FP16 in Pascal (2017) and BF16 in Ampere (2020)), consistently yields twofold or fourfold increases in computational throughput with unchanged energy consumption and silicon footprint. We survey Hopper's advancements in these aspects in \cref{subsec:hardware}.

A significant technical hurdle involves reengineering \textsc{FlashAttention-2} to effectively leverage these advanced hardware resources: asynchrony entails structuring computation so that matrix multiplication and softmax operations can overlap, despite inherent output dependencies, while low-precision arithmetic necessitates strategies for curbing quantization errors—an especially acute concern when handling outlier activations in large language models~\citep{dettmers2208llm, sun2024massive}.

In pursuit of this goal, we introduce \textsc{FlashAttention-3}, which presents and integrates three innovative strategies aimed at optimizing performance for contemporary GPU platforms:\footnote{While our experimental focus is on NVIDIA's Hopper architecture, the methodology generalizes to any GPU offering adequate asynchronous processing and support for low-precision operations.}

\begin{enumerate}

\item \textbf{Producer-Consumer asynchrony:} We propose a warp-specialized software pipelining approach that capitalizes on the asynchronous capabilities of data movement and Tensor Core operations. By segregating data producers and consumers into distinct warps, this scheme augments the algorithm’s capacity to conceal latencies related to both memory accesses and instruction scheduling.

\item \textbf{Hiding softmax under asynchronous block-wise GEMMs:} Our strategy overlaps the relatively slower non-GEMM computations required by softmax, including floating point multiply-add and exponential operations, with the asynchronous execution of WGMMA GEMM instructions. To facilitate this, we reformulate portions of the \textsc{FlashAttention-2} algorithm, thereby eliminating certain intrinsic dependencies between softmax processing and GEMMs. For instance, in the two-stage implementation, while softmax operates on one block of the scores matrix, the asynchronous proxy employing WGMMA concurrently computes the subsequent block.

\item \textbf{Hardware-accelerated low-precision GEMM:} The forward pass algorithm is modified to support execution on FP8 Tensor Cores for GEMM workloads, resulting in an almost twofold increase in observed TFLOPs/s. Accommodating this enhancement requires resolving disparities in WGMMA’s layout assumptions for FP32 accumulators and FP8 operands in memory. By leveraging block quantization and incoherent processing strategies, we effectively alleviate accuracy degradation associated with transitioning to FP8 precision.
\end{enumerate}

For empirical validation, we assess the performance of \textsc{FlashAttention-3} on the H100 SXM5 GPU across a diverse set of configurations. Our results indicate that (1) in the forward pass, FP16 provides a 1.5-2.0$\times$ acceleration over \textsc{FlashAttention-2} (achieving peak throughput up to 740 TFLOPs/s), and yields a 1.5-1.75$\times$ improvement during the backward pass; (2) FP8 attains nearly 1.2 PFLOPs/s; and (3) at higher sequence lengths, FP16 demonstrates superior throughput while FP8 remains competitive with, and in some setups exceeds, the performance of the leading attention routine from NVIDIA's cuDNN library\footnote{In greater detail: for head dimension 64, \textsc{FlashAttention-3} FP8 outperforms, whereas for head dimensions 128 and 256 it matches the baseline for non-causal cases and trails slightly with causal masking.}. Additionally, we confirm that the FP16 implementation in \textsc{FlashAttention-3} maintains numerical accuracy on par with \textsc{FlashAttention-2}, and improves upon the typical attention computation because all intermediate steps (such as softmax scaling) are stored in FP32. When evaluating FP8 \textsc{FlashAttention-3} employing block quantization and incoherent processing, it proves to be 2.6$\times$ more precise than conventional attention with per-tensor quantization, especially under scenarios with outlier feature values.

We provide open-source access to \textsc{FlashAttention-3} under a permissive license\footnote{\textsc{FlashAttention-3} is available at \texttt{https://github.com/Dao-AILab/flash-attention}}, and we aim for seamless integration with the PyTorch and Hugging Face ecosystems to maximize accessibility for researchers and practitioners.

\section{Background: Multi-Head Attention and GPU Characteristics}
\label{sec:background}

\subsection{Multi-Head Attention}
\label{subsec:multi_head_attn}

Consider input sequences for a single attention head: $\mathbf{Q}, \mathbf{K}, \mathbf{V} \in \mathbb{R}^{N \times d}$, corresponding to queries, keys, and values, where $N$ denotes the sequence length and $d$ represents the dimension per head. The attention mechanism produces output $\mathbf{O}$ as follows:

\begin{equation*}
  \mathbf{S} = \alpha \mathbf{Q} \mathbf{K}^\top \in \mathbb{R}^{N \times N}, \quad \mathbf{P} = \mathrm{softmax}(\mathbf{S}) \in \mathbb{R}^{N \times N}, \quad \mathbf{O} = \mathbf{P}\mathbf{V} \in \mathbb{R}^{N \times d},
\end{equation*}

Here, $\mathrm{softmax}$ is performed along each row, and it is common to use $\alpha = 1/\sqrt{d}$ as the scale parameter. To mitigate potential numerical issues caused by exponentiation, one generally subtracts $\mathrm{rowmax}(\mathbf{S})$ from $\mathbf{S}$. In the context of multi-head attention (MHA), distinct projection matrices for queries, keys, and values are used for each head. This process is efficiently parallelized across different heads and batches to generate the complete output tensor.

Let us consider a scalar loss function $\phi$ and denote by $\mathbf{d}(-) = \partial \phi / \partial (-)$ the corresponding gradient operator. Assuming the presence of the output gradient $\mathbf{dO} \in \mathbb{R}^{N \times d}$, the gradients with respect to $\mathbf{Q}$, $\mathbf{K}$, and $\mathbf{V}$—denoted as $\mathbf{dQ}$, $\mathbf{dK}$, and $\mathbf{dV}$—are computed by applying the chain rule as detailed below:

\begin{align*}
  \mathbf{dV} &= \mathbf{P}^\top \mathbf{dO} \in \mathbb{R}^{N \times d} \\
  \mathbf{dP} &= \mathbf{dO} \mathbf{V}^\top \in \mathbb{R}^{N \times N} \\
  \mathbf{dS} &= \mathrm{dsoftmax}(\mathbf{dP}) \in \mathbb{R}^{N \times N} \\
  \mathbf{dQ} &= \alpha \mathbf{dS} \mathbf{K} \in \mathbb{R}^{N \times d} \\
  \mathbf{dK} &= \alpha \mathbf{dS}^\top \mathbf{Q} \in \mathbb{R}^{N \times d},
\end{align*}

For clarification, the derivative of the softmax with respect to its input is given as $\mathbf{d}s = (\mathrm{diag}(p) - p p^\top)\mathbf{d}p$ for a vector $p = \mathrm{softmax}(s)$; we use the notation $\mathrm{dsoftmax}(\mathbf{dP})$ to signify the application of this derivative row-wise. These gradient calculations, in turn, can be efficiently computed in parallel over both the attention heads and the batch dimension during the backward propagation phase of Multi-Head Attention (MHA).

\subsection{GPU hardware characteristics and execution model}
\label{subsec:hardware}

We describe the aspects of the GPU's execution model relevant for \textsc{FlashAttention-3}, with a focus on the NVIDIA Hopper architecture as a concrete instantiation of this model.

\paragraph{Memory hierarchy:} The memory architecture of a GPU is structured as a hierarchy of distinct data domains, where higher bandwidth comes at the cost of reduced capacity (\cref{tab:gpu-hierarchy})\footnote{\citet{luo2024benchmarking} provides a shared memory bandwidth value of 128 bytes per clock cycle per SM; this figure is then scaled by 132 SMs and the maximum clock speed of 1830 MHz.}. At the top is global memory (GMEM), also referred to as HBM, which comprises the off-chip DRAM accessible from any streaming multiprocessor (SM). Data residing in GMEM is automatically cached within a dedicated on-chip L2 cache. Within each SM, there exists a compact, highly banked cache, managed directly by the programmer, known as shared memory (SMEM). Finally, at the most granular level, each SM is equipped with its own register file.

\paragraph{Thread hierarchy:} In the GPU programming paradigm, execution units are systematically arranged into logical collections known as threads. This hierarchy, ordered from the most granular to the broadest, consists of individual threads, followed by warps (each consisting of 32 threads), warpgroups (which group together 4 consecutive warps), then threadblocks (also referred to as cooperative thread arrays, or CTAs), advancing to threadblock clusters (specific to Hopper), and, finally, grids.

The relationship between these two hierarchies is tightly coupled. Threads belonging to the same CTA are scheduled together onto the same SM, while CTAs grouped within the same cluster are simultaneously assigned to the same GPC. SMEM is accessible to every thread inside a CTA, in contrast to RMEM, where each thread possesses up to 256 private registers.

\paragraph{Asynchrony and warp-specialization:}

GPUs achieve high throughput by leveraging parallelism and asynchronous execution to mask delays in memory access and computation. Hopper introduces the Tensor Memory Accelerator (TMA), a specialized hardware module, to enable asynchronous data transfers between GMEM and SMEM \cite[\S7.29]{cuda}. Additionally, Hopper's Tensor Core supports asynchronous operation via the warpgroup-scoped WGMMA instruction \cite[\S9.7.14]{ptx}, allowing it to fetch operands directly from shared memory—an improvement over earlier architectures like Ampere.

Hardware mechanisms that enable asynchrony make it possible to implement warp-specialized kernels, wherein the warps within a CTA are assigned exclusively to either producer (handling data transfers) or consumer (executing computations) tasks. Such division enhances the compiler’s capacity to produce superior instruction schedules \citep{warp-specialization-2011}. Moreover, Hopper provides functionality for dynamically adjusting register allocation among warpgroups through the \verb|setmaxnreg| command \cite[\S9.7.17.1]{ptx}, permitting warps responsible for MMAs to access a greater portion of RMEM relative to those limited to issuing TMAs (which require only one active thread).

\paragraph{Low-precision number formats:}
\label{sec:low-precision-gpu}
Recent GPUs are equipped with dedicated hardware designed to expedite computations using low-precision formats. As an illustration, the WGMMA operation is able to utilize the FP8 Tensor Cores on Hopper, achieving double the throughput per SM compared to FP16 or BF16.

Nevertheless, effectively utilizing FP8 WGMMA requires familiarity with the operand layout requirements. Consider the GEMM operation $A \times B^{\top}$, where $A$ is an $M\times K$ matrix and $B$ is an $N\times K$ matrix. The operand $A$ (or $B$) is termed \emph{mn-major} when it has contiguous elements along the outer $M$ (or $N$) axis, whereas it is designated \emph{k-major} if contiguity is along the inner $K$ dimension. For FP16 WGMMA instructions, input operands in SMEM can be either mn-major or k-major. In contrast, FP8 WGMMA restricts acceptable operand formats to k-major only. This restriction becomes particularly problematic in scenarios like attention mechanisms, where combining consecutive GEMMs within a single kernel is desirable; incompatibilities between the layout of FP32 accumulators and FP8 operands complicate issuing sequential dependent FP8 WGMMAs.

In the context of attention, these layout restrictions entail certain modifications to the design of an FP8 algorithm, which we describe in \cref{sec:algofp8}.

\subsection{Standard Attention and Flash Attention} In accordance with the terminology of \citet{dao2022flashattention}, we use \textbf{standard attention} to refer to the conventional approach for computing attention on the GPU, where the intermediate matrices $\mathbf{S}$ and $\mathbf{P}$ are stored in HBM. The central innovation of \textsc{FlashAttention} is the utilization of a local version of the softmax operation to significantly reduce costly intermediate memory accesses, thereby consolidating the entire attention computation into a single kernel. This local softmax is performed in lines \ref{code-ws:softmax_start}-\ref{code-ws:softmax_end} within the consumer mainloop of \cref{alg:flash3_wgmma_ws_only}, along with appropriate rescaling of blocks of $\mathbf{O}$. A straightforward proof demonstrating that this method correctly produces $\mathbf{O}$ is provided in \cite[\S 2.3.1]{dao2023flashattention2}.

\section{FlashAttention-3: Algorithm}
\label{sec:algo}

This section outlines the \textsc{FlashAttention-3} algorithm. For clarity, the discussion centers on the forward computation, while details of the backward pass are presented in~\cref{sec:algo_ws_bwd}. Initially, we demonstrate the incorporation of warp specialization alongside a circular SMEM buffer into the foundational framework of \textsc{FlashAttention-2}. Next, we detail the manner in which WGMMA's asynchronous operations facilitate the construction of a two-phase pipeline that interleaves GEMM and softmax computations. Lastly, we discuss the necessary adaptations to enable FP8 support, addressing both format compliance and precision through block quantization and non-uniform processing.

\subsection{Producer-Consumer asynchrony through warp-specialization and pingpong scheduling}
\label{sec:algo_ws}

\paragraph{Warp-specialization}
Similar to \textsc{FlashAttention-2}, the forward pass in \textsc{FlashAttention-3} features extensive parallelism along the batch dimension, the number of heads, and the length of the query sequence. Accordingly, a conceptual explanation at the CTA-level—processing a tile $\mathbf{Q}_i$ of the input query matrix and generating the respective output tile $\mathbf{O}_i$—is adequate. For clarity, we begin by outlining the warp-specialization approach utilizing a circular SMEM buffer, initially \emph{excluding} the GEMM-softmax overlapping. Denote by $d$ the dimension of each attention head, $N$ the sequence length, and select a block size $B_r$ for queries, allowing partition of $\mathbf{Q}$ into $T_r = \lceil \frac{N}{B_r} \rceil$ blocks, namely $\mathbf{Q}_1, .., \mathbf{Q}_{T_r}$.

\begin{algorithm}[H]
    \caption{\small\label{alg:flash3_wgmma_ws_only}\textsc{FlashAttention-3} forward pass \textbf{with no} intra-consumer overlap -- CTA perspective}
    \begin{algorithmic}[1]
\REQUIRE Input matrices $\mathbf{Q}_i \in \mathbb{R}^{B_r \times d}$, $\mathbf{K}, \mathbf{V} \in \mathbb{R}^{N \times d}$ residing in HBM; key block size $B_c$ with $T_c = \lceil \frac{N}{B_c} \rceil$.
\STATE Create a pipeline handler to orchestrate barrier synchronization, utilizing an $s$-stage circular SMEM buffer.
\IF {executing within producer warpgroup}
\STATE Release a pre-specified count of registers.
\STATE Begin by loading $\mathbf{Q}_i$ from HBM into shared memory.
\STATE Once this operation completes, signal consumers that $\mathbf{Q}_i$ is now available in shared memory.
\FOR{$0 \le j < T_c$}
    \STATE Await consumption of the $(j\,\%\,s)$ buffer stage before proceeding.
    \STATE Initiate the transfer of $\mathbf{K}_j$ and $\mathbf{V}_j$ from HBM to the $(j\,\%\,s)$ slot in shared memory.
    \STATE After the load is done, inform consumers that $\mathbf{K}_j$ and $\mathbf{V}_j$ are now present.
\ENDFOR
\ELSE \STATE Acquire the appropriate number of registers, determined by how many consumer warps exist.
\STATE Allocate $\mathbf{O}_i = (0) \in \mathbb{R}^{B_r \times d}$ and initialize $\ell_i, m_i = (0), (-\infty) \in \mathbb{R}^{B_r}$ on-chip.
\STATE Wait for confirmation that $\mathbf{Q}_i$ is present in shared memory.
\FOR{$0 \le j < T_c$}
\STATE Await the completion of $\mathbf{K}_j$ loading into shared memory.
\STATE Perform the matrix multiplication $\mathbf{S}_i^{(j)} = \mathbf{Q}_i \mathbf{K}_j^T$ (SS-GEMM). Commit the operation and wait. \label{alg:ws_only_gemm1}
\STATE Save $m_i^{\mathrm{old}} = m_i$, then update $m_i = \mathrm{max}(m_i^{\mathrm{old}}, \mathrm{rowmax}(\mathbf{S}_i^{(j)}))$. \label{code-ws:softmax_start}
\STATE Calculate $\widetilde{\mathbf{P}}_i^{(j)} = \mathrm{exp}(\mathbf{S}_i^{(j)} - m_i)$ and refresh $\ell_i = \mathrm{exp}(m_i^{\mathrm{old}} - m_i) \ell_i + \mathrm{rowsum}(\widetilde{\mathbf{P}}_i^{(j)})$. \label{code-ws:softmax_end}
\STATE Wait for $\mathbf{V}_j$ to become available in shared memory.
\STATE Update $\mathbf{O}_i = \mathrm{diag}(\mathrm{exp}(m_i^{\mathrm{old}} - m_i))^{-1} \mathbf{O}_i + \widetilde{\mathbf{P}}_i^{(j)} \mathbf{V}_j$ (RS-GEMM). Commit this operation and wait for it to complete. \label{alg:ws_only_gemm2}
\STATE Free up the $(j\,\%\,s)$ buffer stage so the producer can continue.
\ENDFOR
\STATE Finalize by evaluating $\mathbf{O}_i = \mathrm{diag}(\ell_i)^{-1} \mathbf{O}_i$ and $L_i = m_i + \log(\ell_i)$.
\STATE Store results $\mathbf{O}_i$ and $L_i$ to HBM, writing them as the $i$th partitions of $\mathbf{O}$ and $L$, respectively.
\ENDIF
\end{algorithmic}
\end{algorithm}

In our deployment of \cref{alg:flash3_wgmma_ws_only} on the Hopper architecture, we employ \verb|setmaxnreg| for handling (de)allocations, utilize TMA for fetching $\mathbf{Q}_i$ as well as $\{ \mathbf{K}_j, \mathbf{V}_j \}_{0 \leq j < T_c}$, and leverage WGMMA to perform GEMMs within the consumer mainloop. The SS and RS prefixes indicate whether the initial operand is obtained from shared memory or the register file, respectively. When analyzing the execution pattern of \cref{alg:flash3_wgmma_ws_only}, it is important to understand that TMA loads are asynchronous and thus initiating new loads will not be blocked by the completion status of previous ones. Additionally, within the producer mainloop, no wait instructions are triggered during the initial $s$ cycles since the buffer is still being populated.

\paragraph{Pingpong scheduling} The decoupled execution provided by the asynchronous behavior of WGMMA and TMA, combined with warp-specialization, allows for softmax computations in one warpgroup to proceed simultaneously with GEMM computations in a different warpgroup. The motivation for this scheduling technique arises from the significant disparity in computational throughput between matmul and non-matmul operations on contemporary hardware accelerators. For instance, an H100 SXM5 GPU supports up to 989 TFLOPS for FP16 matmul computations, while its capacity for executing special function operations such as the exponential\footnote{According to the CUDA programming guide, each streaming multiprocessor (SM) can execute 16 special function operations per clock. Multiplying 16 by the 132 SMs and a 1830 MHz clock frequency produces an estimated 3.9 TFLOPS for special functions.} (a core component of softmax) is restricted to only 3.9 TFLOPS. Considering the attention forward pass with FP16 and a head size of 128, the matmul component demands 512 times as many floating-point operations as the exponential, even though the latter operates at just $1/256$th the throughput. This disparity means the time spent on the exponential operation can consume as much as 50\% of the cycles needed for matmul. The imbalance is exacerbated in the FP8 format, which doubles the matmul throughput without increasing the throughput for exponentials.

Because the exponential computation is handled by a dedicated hardware unit (the multi-function unit), optimal performance is achieved when the exponential operation overlaps with matmul executions on the Tensor Cores. To arrange this, we utilize synchronization barriers (\texttt{bar.sync} instructions) to control the execution order: the GEMMs (GEMM1 -- $\mathbf{P} \mathbf{V}$ from one step, and GEMM0 -- $\mathbf{Q} \mathbf{K}^\top$
from the next step) associated with warpgroup 1 are made to run prior to the GEMMs for warpgroup 2. This setup ensures that warpgroup 1's softmax computation is active concurrently with the GEMMs of warpgroup 2. The process then alternates, enabling warpgroup 2 to compute softmax while warpgroup 1 takes up GEMMs; this staggered execution is referred to as ``pingpong'' scheduling, as depicted in~\cref{fig:pingpong_scheduling}. Although in practice the alternation is less ideal than the illustration might suggest, this technique typically boosts throughput (for example, from 570 TFLOPS to between 620 and 640 TFLOPS in the FP16 forward pass, using a head size of 128 and a sequence length of 8192).

\paragraph{Attention variants} In the context of multi-query attention~\citep{shazeer2019fast} and grouped query attention~\citep{ainslie2023gqa}, we adopt the strategy from \textsc{FlashAttention-2}, where tensor indexing is adapted so as to prevent unnecessary replication of $\mathbf{K}$ and $\mathbf{V}$ tensors in HBM.

\subsection{Intra-warpgroup overlapping GEMMs and softmax}

Even within one warpgroup, we can overlap some instructions in the softmax with some instructions in the GEMMs. We describe one technique to do so.

Within the attention algorithm, the primary (inner) loop exhibits inherent sequential dependencies that limit opportunities for parallel execution within a single iteration. Specifically, the (local) softmax operation (see lines \ref{code-ws:softmax_start} through \ref{code-ws:softmax_end}) is dependent on the output $\mathbf{S}_i^{(j)}$ produced by the initial GEMM, while the subsequent GEMM consumes the intermediate result $\widetilde{\mathbf{P}}_i^{(j)}$. These dependencies are enforced by explicit wait statements at lines \ref{alg:ws_only_gemm1} and \ref{alg:ws_only_gemm2} in \cref{alg:flash3_wgmma_ws_only}, thereby serializing the processing of softmax and GEMM steps. Nevertheless, by introducing extra register buffers and employing pipelining across loop iterations, these dependencies can be alleviated. Building on this strategy, we introduce a two-stage\footnote{The number of pipeline stages can be up to, but does not necessarily have to match, the number $s$ of stages present in the circular SMEM buffer.} GEMM-softmax pipelining method:

\begin{algorithm}[H]
    \caption{\small\label{alg:flash3_wgmma}\textsc{FlashAttention-3} consumer warpgroup forward pass}
    \begin{algorithmic}[1]
      \REQUIRE  Input matrices $\mathbf{Q}_i \in \mathbb{R}^{B_r \times d}$ and $\mathbf{K}, \mathbf{V} \in \mathbb{R}^{N \times d}$ located in HBM, specified key block size $B_c$ so that $T_c = \lceil \frac{N}{B_c} \rceil$.
      \STATE Dynamically assign a predetermined set of registers depending on the current count of consumer warps.
      \STATE Initialize $\mathbf{O}_i = (0) \in \mathbb{R}^{B_r \times d}$ and set both $\ell_i, m_i$ to $(0), (-\infty) \in \mathbb{R}^{B_r}$ in on-chip memory.
      \STATE Synchronize until $\mathbf{Q}_i$ and $\mathbf{K}_0$ have been fully transferred into shared memory.
      \STATE Evaluate $\mathbf{S}_{\mathrm{cur}} = \mathbf{Q}_i \mathbf{K}_0^T$ via WGMMA operation; commit and synchronize.
      \STATE Free the buffer's initial stage allocated for $\mathbf{K}$.
      \STATE Derive $m_{i}$, $\tilde{\mathbf{P}}_{\mathrm{cur}}$, and $\ell_{i}$ from $\mathbf{S}_{\mathrm{cur}}$; apply rescaling to $\mathbf{O}_i$.
      \FOR{$1 \le j < T_c - 1$}
        \STATE Wait until $\mathbf{K}_j$ is present in shared memory.
        \label{code:mainloop_start}
        \STATE Compute $\mathbf{S}_{\mathrm{next}} = \mathbf{Q}_i \mathbf{K}_{j}^T$ with WGMMA. Commit this operation, do not synchronize yet. \label{code:first_wgmma}
        \STATE Block until $\mathbf{V}_{j-1}$ has loaded into shared memory.
        \STATE Update $\mathbf{O}_{i}$ according to $\mathbf{O}_{i} = \mathbf{O}_{i} + \tilde{\mathbf{P}}_{\mathrm{cur}} \mathbf{V}_{j-1}$, employing WGMMA. Commit, but delay synchronization. \label{code:second_wgmma}
        \STATE Await completion of the WGMMA task computing $\mathbf{Q}_i \mathbf{K}_{j}^T$.
        \STATE Using $\mathbf{S}_{\mathrm{next}}$, calculate $m_{i}$, $\tilde{\mathbf{P}}_{\mathrm{next}}$, and $\ell_{i}$. \label{code:softmax}
        \STATE Await finalization of the WGMMA $\tilde{\mathbf{P}}_{\mathrm{cur}} \mathbf{V}_{j-1}$, then apply the appropriate rescale to $\mathbf{O}_i$.
        \STATE Release buffer stages: $(j\,\%\,s)$th for $\mathbf{K}$, and $(j-1\,\%\,s)$th for $\mathbf{V}$.
        \STATE Move contents of $\mathbf{S}_{\mathrm{next}}$ into $\mathbf{S}_{\mathrm{cur}}$.
        \label{code:mainloop_end}
      \ENDFOR \STATE Wait for the loading of $\mathbf{V}_{T_c - 1}$ into shared memory.
      \STATE Perform the update $\mathbf{O}_{i} = \mathbf{O}_{i} + \tilde{\mathbf{P}}_{\mathrm{last}} \mathbf{V}_{T_c - 1}$ by WGMMA; commit and wait for completion.

\cref{alg:flash3_wgmma} serves as a substitute for the consumer pathway specified in \cref{alg:flash3_wgmma_ws_only}, thereby forming the complete \textsc{FlashAttention-3} procedure for FP16 arithmetic. Broadly, we refer to WGMMA as a shorthand for asynchronous GEMM. During the mainloop, which spans from lines \ref{code:mainloop_start} to \ref{code:mainloop_end}, the execution of the second WGMMA kernel for iteration $j$ (line \ref{code:second_wgmma}) is orchestrated concurrently with the softmax computation for the next iteration $j+1$ (line \ref{code:softmax}).

Although the pipelined architecture depicted above promises theoretical improvements in throughput, a number of real-world factors must be addressed:
\paragraph{Compiler reordering} The pseudocode is written with an idealized sequence of instructions, but in practice, the NVCC compiler may rearrange the order of operations for optimization purposes. Such instruction reordering can interfere with the intended interleaving of WGMMA and non-WGMMA computations, potentially reducing or nullifying the anticipated performance advantages. SASS code inspection indicates, however, that the generated machine code still overlaps instructions as expected (Section~\ref{sec:2-stage-sass}).
\paragraph{Register pressure} Minimizing register spilling is essential for achieving high performance. The 2-stage pipeline, however, demands additional registers to accommodate intermediate values and manage state between stages. In particular, the algorithm necessitates an extra $\mathbf{S}_{\mathrm{next}}$ held in registers, corresponding to an added register allocation of $B_r \times B_c \times \text{sizeof}(\text{float})$ per threadblock. This greater register requirement can compete with other optimizations such as increasing block sizes, as both approaches consume more registers. In practice, profiling should guide the balance between these competing demands.
\paragraph{3-stage pipelining} Building on the prior 2-stage scheme, we develop a 3-stage version designed to overlap the second WGMMA operation with the softmax phase. While this deeper pipeline could provide further improvements in Tensor Core throughput, it entails even higher register consumption due to the additional pipeline stage, thus intensifying the complexity of managing the trade-off between pipeline depth and tile dimensions. Section~\cref{sec:3-stage} contains a comprehensive explanation of the 3-stage approach alongside its empirical assessment.

\subsection{Low-precision with FP8}
\label{sec:algofp8}

\textbf{Efficiency: layout transformations.} Executing the forward computation of \textsc{FlashAttention-3} at FP8 precision introduces unique complications related to memory layout compatibility that are not present when using FP16.

Typically, the input tensors $\mathbf{Q}$, $\mathbf{K}$, and $\mathbf{V}$ are stored such that they are contiguous along the head dimension. However, for FP8 WGMMA’s k-major format required by the second GEMM, it is necessary that $\mathbf{V}$—specifically, the sections of $\mathbf{V}$ loaded into SMEM—be contiguous in the sequence length dimension. Since TMA loads do not modify which dimension is contiguous, this necessitates either (1) transposing $\mathbf{V}$ in global memory prior to computation, or (2) performing an in-kernel transposition of the tiles of $\mathbf{V}$ after transferring them into SMEM. For strategy (1), one can (1a) combine the transpose operation with the epilogue phase of a preceding computation such as the rotary embedding, or (1b) utilize a dedicated pre-processing transpose kernel\footnote{A well-optimized transpose kernel can achieve performance close to the theoretical memory bandwidth of the device~\citep{colfax_cutlass_transpose_2024}.} to switch the strides between the sequence length and head dimensions. Nevertheless, (1a) presents integration challenges within standard libraries, while (1b) incurs excessive cost in memory bandwidth—making it undesirable, particularly for memory-bound scenarios like inference.

For FP8 \textsc{FlashAttention-3}, we choose the second approach. During the in-kernel transpose, we exploit the capabilities of the LDSM (\verb|ldmatrix|) and STSM (\verb|stmatrix|) instructions. These instructions enable a warp of threads to collaboratively transfer data between shared memory (SMEM) and registers (RMEM), working on chunks of 128 bytes at a time.\footnote{According to the PTX documentation, LDSM and STSM are formally specified as moving $8 \times 8$ matrices composed of 16-bit elements \cite[\S 9.7.13.4.15-16]{ptx}; nevertheless, by pairing two 8-bit values, we can adapt these instructions to suit FP8-level operations. On the other hand, the transpose-specific LDSM/STSM instructions cannot independently handle such packed 8-bit data, which requires additional register manipulations between LDSM and STSM to properly accomplish tile-level transpositions; the fine-grained mechanics are excluded here.} These LDSM/STSM operations are economical in terms of register usage, permitting their invocation within the producer warpgroup, and can perform transposition concurrently with data movement. Additionally, after the initial pass, the subsequent transposition of the next $\mathbf{V}$ tile can be overlapped with the dual WGMMAs that process the previous $\mathbf{V}$ and the current $\mathbf{K}$ tile, further optimizing computation overlap.

Secondly, we note that, in contrast to FP16, the arrangement of the FP32 accumulator in an FP8 WGMMA does not align with the register layout of operand A. Illustrative excerpts from these layouts are provided in \cref{fig:rmem_accum} and \cref{fig:rmem_operand}, where register entries are assigned per thread according to the order shown. Through the application of byte permute instructions, the accumulator output of the initial WGMMA can be reorganized to fit the required format for a subsequent WGMMA call and to align with the memory layout of the $\mathbf{V}$ tile generated by the in-kernel transpose. More precisely, as depicted in \cref{fig:rmem_accum}, the register sequence is reordered as
$$\{ \verb|d0 d1 d4 d5 d2 d3 d6 d7| \},$$
with this arrangement repeated every 8 bytes. Conceptually, this operation reorganizes the columns within the $\mathbf{P}$ tile—for example, columns $0189$ now occupy the initial four positions. To ensure WGMMA processes the correct output tile, the in-kernel transpose can accordingly be modified to emit a correspondingly permuted sequence of rows for the $\mathbf{V}$ tile.\footnote{By incorporating this flexible approach through an in-kernel transpose, there is no longer a need for shuffle instructions to reassign registers among threads, as previously discussed in~\citep{colfax_fp8_flashattention_2024}.}

\textbf{Accuracy: block quantization and incoherent processing.} The FP8 (e4m3) data type allocates only 3 bits to represent the mantissa and 4 bits for the exponent, which introduces a greater degree of numerical inaccuracy relative to FP16/BF16 formats. The challenge is further exacerbated in large-scale models, as outlier values can occur~\citep{dettmers2208llm,
  sun2024massive} with magnitudes that vastly exceed the majority of elements, thereby complicating the quantization process. Standard practice involves per-tensor scaling~\citep{micikevicius2022fp8}, whereby a single scale factor is stored for each tensor (such as one for each of $\mathbf{Q}$, $\mathbf{K}$, and $\mathbf{V}$).
In order to mitigate the quantization errors in FP8-based attention, we incorporate two primary strategies:

\begin{enumerate}

\item \textbf{Block quantization}: Instead of a global scale, we retain a distinct scaling factor for each block, segmenting $\mathbf{Q}$, $\mathbf{K}$, and $\mathbf{V}$ tensors into chunks of dimension $B_r \times d$ or $B_c \times d$ and applying quantization independently to each. This block-wise quantization can be seamlessly integrated with preceding operations, such as rotary embedding, without incurring additional latency—owing to the memory bandwidth limitations of rotary embedding. As \textsc{FlashAttention-3} works intrinsically on blocks, the block-wise scaling of $\mathbf{S}$ can be performed concurrently with no overhead in computation.

\item \textbf{Incoherent processing}: To mitigate the effects of outliers, $\mathbf{Q}$ and $\mathbf{K}$ are each right-multiplied by a random orthogonal matrix $\mathbf{M}$ prior to being quantized to FP8 format. Given that $\mathbf{M}$ is orthogonal (so $\mathbf{M}\mathbf{M}^\top = I$), it follows that $(\mathbf{Q}\mathbf{M})(\mathbf{K}\mathbf{M})^\top = \mathbf{Q}\mathbf{K}^\top$, implying this operation leaves the attention computation unchanged. The application of $\mathbf{M}$ has a 'spreading' effect, redistributing outlier values since every element in $\mathbf{Q}\mathbf{M}$ or $\mathbf{K}\mathbf{M}$ constitutes a random linear combination of the originals, resulting in lower quantization error. Consistent with the methodology proposed by \citet{chee2024quip} and \citet{tseng2024quip}, $\mathbf{M}$ is constructed as a product of random diagonal matrices of entries $\pm 1$ and a Hadamard matrix; this particular structure allows for matrix-vector multiplication in $O(d \log d)$ time (rather than $O(d^2)$), and, similarly to rotary embedding, can be fused at no additional computational expense.
\end{enumerate}

We validate that these two techniques reduces numerical error by up to 2.6$\times$ in \cref{sec:numerical_error}.

\section{Empirical Validation}
\label{sec:exp}

We use the primitives from CUTLASS~\citep{Thakkar_CUTLASS_2023} such as WGMMA and TMA abstractions to implement \textsc{FlashAttention-3} and evaluate its efficiency and accuracy.

\begin{itemize}

\item \textbf{Attention Benchmarking.} We evaluate the performance of \textsc{FlashAttention-3} over a range of sequence lengths, comparing its execution time with that of the conventional PyTorch implementation, \textsc{FlashAttention-2}, a version of \textsc{FlashAttention-2} implemented in Triton (leveraging H100-specific instructions), and a cuDNN-based, vendor-optimized version of \textsc{FlashAttention-2} targeting H100 GPUs. The results indicate that \textsc{FlashAttention-3} achieves speedups of up to 2.0$\times$ relative to \textsc{FlashAttention-2} and is 1.5$\times$ faster compared to the Triton variant of \textsc{FlashAttention-2}. Furthermore, \textsc{FlashAttention-3} attains throughput up to 740 TFLOPs/s, amounting to 75\% of the H100 GPUs’ peak theoretical TFLOPs/s.

\item \textbf{Ablation Study.} Our experimental ablations show that the incorporation of warp-specialization and the pipelining of GEMM with softmax operations are key contributors to the improved performance observed in \textsc{FlashAttention-3}.

\item \textbf{FP8 Attention Accuracy.} We demonstrate that the adoption of block quantization together with incoherent processing techniques lowers the numerical errors of FP8 \textsc{FlashAttention-3} by a factor of 2.6$\times$.

\subsection{Benchmarking Attention}
\label{subsec:benchmark_attn}

We evaluate the execution times of various attention algorithms using an H100 80GB SXM5 GPU, considering both the presence and absence of a causal mask and head dimensions of either 64 or 128 for FP16 input data. The outcomes, presented in~\cref{fig:benchmark_attn_fwd} and~\cref{fig:benchmark_attn_bwd}, demonstrate that \textsc{FlashAttention-3} achieves speedups of approximately 1.5-2.0$\times$ over \textsc{FlashAttention-2} during the forward computation and about 1.5-1.75$\times$ faster performance in the backward computation. When compared against a traditional attention approach, \textsc{FlashAttention-3} exhibits acceleration by factors ranging from 3$\times$ to 16$\times$. Notably, for sequence lengths of 1k or more, \textsc{FlashAttention-3} outperforms even the cuDNN library (a proprietary solution from the hardware vendor) which is specifically tuned for the H100 architecture.

\paragraph{Benchmark settings:} We experiment with sequence lengths of 512, 1k, up to 16k, choosing the batch size such that the aggregate token count is fixed at 16k. The hidden size is set to 2048, while the head dimension is selected from 64, 128, or 256, corresponding to 32, 16, or 8 attention heads, respectively. The number of FLOPs for the forward pass is determined using:

\begin{equation*}
  4 \cdot \text{seqlen}^2 \cdot \text{head dimension} \cdot \text{number of heads}.
\end{equation*}

When applying causal masking, this value is halved, reflecting that only roughly 50\% of the elements are actually computed. For estimating the FLOPs required by the backward pass, the FLOPs from the forward computation are scaled by a factor of 2.5. This ratio arises from the structure of the computation: the forward pass involves 2 matrix multiplications, whereas the backward pass (which employs recomputation) comprises 5 matrix multiplications.

Additionally, we evaluate and record the forward pass runtime when using FP8 precision, keeping experimental conditions consistent. The results pertaining to headdim 256 are presented in ~\cref{fig:benchmark_attn_fp8}, while comprehensive data can be found in \cref{sec:benchmark_attn_fp8_full}.

\subsection{Ablation Study: 2-Stage Pipelining Experiments}
\label{sec:2-stage-pipelining-experiments}

We conduct an ablation study on the 2-stage WGMMA-softmax pipelining and the warp-specialization technique in the context of non-causal FP16 \textsc{FlashAttention-3}, using fixed values $\{\text{batch}, \text{seqlen}, \text{nheads}, \text{hdim}\} = \{ 4, 8448, 16, 128\}$. As presented in~\cref{table:ablation_pipelining}, these algorithmic optimizations—introducing asynchrony via warp-specialization and enabling the overlap of GEMM and softmax—produce notable performance gains, improving throughput from 570 to 661 TFLOPs.

\subsection{Numerical Error Validation}
\label{sec:numerical_error}

Given the recent focus on the numerical inaccuracies associated with \textsc{FlashAttention}~\citep{golden2024flash}, we evaluate and contrast the performance of \textsc{FlashAttention-2}, \textsc{FlashAttention-3}, and a conventional attention implementation against a high-precision FP64 reference. To model the presence of extreme values in LLM features and activations~\citep{dettmers2208llm, sun2024massive}, we sample the elements of $\mathbf{Q}, \mathbf{K}, \mathbf{V}$ according to the following probability distribution:

\begin{equation*}
  \mathcal{N}(0, 1) + \mathcal{N}(0, 100) \cdot \mathrm{Bernoulli}(0.001).
\end{equation*}

Specifically, all entries follow a normal distribution with a mean of zero and a standard deviation of 1; however, for 0.1\% of the entries, we add a separate, normally distributed term with a standard deviation of 10. We subsequently compute the root mean squared error (RMSE), as shown in~\cref{table:numerical_error}. Under FP16 precision, both \textsc{FlashAttention-2} and \textsc{FlashAttention-3} yield an RMSE that is 1.7$\times$ lower than that of the conventional implementation, owing to the retention of intermediate (softmax) computations in FP32. For the baseline attention in FP8, per-tensor scaling is applied, the matrix multiplication accumulator operates in FP32, and the softmax intermediates are stored in FP16. Benefiting from the use of block quantization and incoherent processing, \textsc{FlashAttention-3} with FP8 surpasses this baseline in accuracy by a factor of 2.6$\times$.

\section{Dicussion, Limitations, Conclusion}
\label{sec:discussion}

By leveraging innovations in both programming and hardware—specifically asynchrony and low-precision arithmetic—\textsc{FlashAttention-3} delivers significant improvements in both efficiency and reliability over previous methods. Our approach achieves a 1.5-2.0$\times$ increase in attention computation speed relative to \textsc{FlashAttention-2}, while also decreasing the numerical error in FP8 by a factor of 2.6 compared to conventional per-tensor quantization. Nonetheless, certain aspects warrant further investigation: enhancing performance for LLM inference scenarios, incorporating a persistent kernel framework within the FP8 kernel,\footnote{For benchmarking purposes, the FP16 variant of \textsc{FlashAttention-3} employs both a persistent kernel and load-balancing strategy, whereas the FP8 version currently does not. This partially accounts for FP8 \textsc{FlashAttention-3}’s weaker performance on short sequence lengths and causal masking tasks relative to FP8 cuDNN kernels.} and systematically exploring the influence of low-precision attention throughout large-scale model training. Although this study centers on Hopper GPUs, the presented strategies are anticipated to transfer effectively to a broader range of hardware accelerators. We envisage that more efficient and accurate attention mechanisms, such as the one introduced here, will facilitate novel use cases—particularly in tasks requiring extended contexts.

\end{document}